\section{Conclusion}\label{sec:conclusion}

This paper documents the distributional profile of out-of-pocket health expenditure in Peru using conditional quantile regression on the 2024 ENAHO. Three findings emerge.

First, the correlates of OOP health expenditure share are quantile-dependent. OLS estimates that SIS membership has no significant average association with OOP share mask heterogeneity across the conditional distribution: at $\tau = 0.75$ and $\tau = 0.90$, SIS--quintile interactions indicate that insured households in higher consumption quintiles have lower conditional OOP shares than their uninsured counterparts. This quantile heterogeneity is invisible in standard CHE analyses based on binary thresholds or conditional-mean regressions.

Second, the SIS--quintile interaction pattern forms a monotonic protection gradient rather than a middle-income squeeze. The magnitude of the SIS association with OOP share increases steadily from Q1 (no significant association) to Q5 ($-1.94$ percentage points at $\tau = 0.90$), and from $\tau = 0.50$ to $\tau = 0.90$. This pattern is more consistent with regressive financial protection---insurance coverage is associated with larger OOP reductions for higher-income enrollees---than with the hypothesis that middle-income households face disproportionate burden relative to both the poorest and the wealthiest. The monotonic gradient is itself a substantive finding with policy implications, even though it does not support the initial framing.

Third, consumption quintile gradients steepen monotonically across the conditional distribution. At $\tau = 0.90$, the Q5 coefficient is 5.16 percentage points, indicating that the financial exposure of high-spending households is substantially understated by average-based measures. The tail of the OOP distribution, not its center, drives catastrophic spending episodes.

These findings carry policy implications, subject to important caveats. Monitoring CHE through binary threshold indicators alone understates the concentration of financial risk in the upper tail. Distributional measures---such as the conditional quantile gap between insured and uninsured households at $\tau = 0.75$ or $\tau = 0.90$---could provide more informative signals for policy targeting. The urban--rural divergence in the SIS association at upper quantiles suggests that the channels through which insurance correlates with OOP burden differ across contexts. The regressive protection pattern raises questions about whether SIS's financial protection value is concentrated among higher-income enrollees who may be accessing coverage through targeting leakage rather than legitimate enrollment.

Our analysis has several limitations, two of which we classify as requiring resolution before the findings can be considered reliable. First, the EsSalud insurance variable ($N = 648$, approximately 2\% coverage vs.\ the expected 30\% nationally) is almost certainly misconstrued due to the household-head-only assignment rule; this misclassification contaminates the uninsured reference group and biases all insurance comparisons. Correcting the variable to use any-household-member coverage is a prerequisite for credible results. Second, consumption quintiles constructed from \texttt{GASHOG2D} include health expenditure, creating a mechanical positive gradient that inflates quintile--OOP correlations and confounds the insurance--quintile interactions. Quintile reconstruction from non-health consumption is required to isolate behavioral from mechanical associations.

Additional limitations include: the 200-replication cluster bootstrap is below the 999 standard for publishable inference, and significance conclusions at upper-tail quantiles ($\tau = 0.75, 0.90$) are provisional; the 44.1\% zero-OOP mass renders standard quantile regression uninformative at lower quantiles, and we did not implement censored quantile regression \citep{Powell1986, Chernozhukov2002} or a full two-part model as the primary specification; no multiple testing correction was applied; the cross-sectional design precludes causal inference; and the primary placebo test is uninformative.

Future work should address these limitations through several extensions: (i) censored quantile regression \citep{Powell1986, Chernozhukov2002} as the primary estimator to handle the zero mass; (ii) quintiles constructed from consumption net of health expenditure to resolve the circularity concern; (iii) 999+ bootstrap replications for reliable upper-tail inference; (iv) correction of the EsSalud variable construction or pooled ENAHO waves to increase EsSalud cell sizes; (v) RIF unconditional quantile regression \citep{Firpo2009} to provide the population-level complement to conditional estimates; and (vi) quasi-experimental designs exploiting SIS enrollment changes or eligibility discontinuities to identify causal effects. The monotonic protection gradient documented here---whether it reflects regressive insurance design, targeting leakage, or differential care-seeking by income---merits investigation with methods that can distinguish among these mechanisms.
